\documentclass[11pt,a4paper]{article}
% preamble.tex — shared setup for human-verification documents
% Compile each document with:  tectonic <File>.tex

\usepackage[T1]{fontenc}
\usepackage[utf8]{inputenc}
\usepackage{lmodern}
\usepackage{amssymb,amsmath,amsthm}
\usepackage{mathtools}
\usepackage[dvipsnames]{xcolor}
\usepackage{listings}
\usepackage{hyperref}
\usepackage[margin=2.5cm]{geometry}
\usepackage{parskip}

% ── Colours ──────────────────────────────────────────────────────────
\definecolor{leanBg}{HTML}{F7F7F0}
\definecolor{haskellBg}{HTML}{F0F4F7}
\definecolor{keyword}{HTML}{7B3294}
\definecolor{comment}{HTML}{5A8C32}
\definecolor{noteFg}{HTML}{B35900}

% ── Lean 4 listings language ─────────────────────────────────────────
% NOTE: Lean listing blocks use ASCII approximations of unicode symbols
% (e.g. -> for arrow, o. for compose).  The literate key substitutes
% these with proper math-mode renderings in the PDF.
\lstdefinelanguage{Lean4}{
  morekeywords={def,theorem,structure,where,namespace,end,import,fun,
                let,have,show,by,exact,intro,cases,rfl,subst,obtain,
                funext,simp,rw,apply,variable,Type,Prop,instance,
                section,open,class,extends,deriving,sorry,noncomputable},
  sensitive=true,
  morecomment=[l]{--},
  morecomment=[n]{/-}{-/},
  morestring=[b]",
}

% ── Lean listing style ───────────────────────────────────────────────
\lstdefinestyle{lean}{
  language=Lean4,
  backgroundcolor=\color{leanBg},
  basicstyle=\small\ttfamily,
  keywordstyle=\color{keyword}\bfseries,
  commentstyle=\color{comment}\itshape,
  breaklines=true,
  frame=single,
  framerule=0.4pt,
  xleftmargin=4pt,
  xrightmargin=4pt,
  aboveskip=6pt,
  belowskip=4pt,
  literate=
    {->}{{$\to$}}2
    {=>}{{$\Rightarrow$}}2
    {<|}{{$\langle$}}1
    {|>}{{$\rangle$}}1
    {==>}{{$\Longrightarrow$}}3
    { o. }{{$\;\circ\;$}}3
    {forall}{{$\forall$}}1
    {:=}{{$\,:=\,$}}2
    {!=}{{$\neq$}}2
    {<=}{{$\leq$}}2
    {>=}{{$\geq$}}2,
}

% ── Haskell listing style ───────────────────────────────────────────
\lstdefinestyle{haskell}{
  language=Haskell,
  backgroundcolor=\color{haskellBg},
  basicstyle=\small\ttfamily,
  keywordstyle=\color{keyword}\bfseries,
  commentstyle=\color{comment}\itshape,
  breaklines=true,
  frame=single,
  framerule=0.4pt,
  xleftmargin=4pt,
  xrightmargin=4pt,
  aboveskip=6pt,
  belowskip=4pt,
  literate=
    {->}{{$\to$}}2
    {=>}{{$\Rightarrow$}}2,
}

% ── Theorem environments ────────────────────────────────────────────
\theoremstyle{definition}
\newtheorem{defn}{Definition}[section]
\newtheorem{thm}[defn]{Theorem}
\newtheorem{prop}[defn]{Proposition}

\theoremstyle{remark}
\newtheorem*{fidelitynote}{Fidelity note}

% ── Convenience macros ───────────────────────────────────────────────
\newcommand{\lean}[1]{\texttt{#1}}      % inline Lean identifier
\newcommand{\hask}[1]{\texttt{#1}}      % inline Haskell identifier
\newcommand{\Ext}{\mathrm{Ext}}
\newcommand{\Pos}{\mathrm{Pos}}
\newcommand{\Shape}{\mathrm{Shape}}


\title{Human Verification: \texttt{Containers.Cont}\\[4pt]
  \large The Category of Containers: Identity, Composition, and Finite (Co)Products}
\author{Machine-generated from \texttt{lean/Cont.lean} for review by Neil Ghani}
\date{July 2026}

\begin{document}
\maketitle

\begin{abstract}
This document presents every definition and theorem \emph{statement} from the
Lean~4 module \lean{Containers.Cont} in three parallel forms: the original
Lean, a readable Haskell rendering, and precise mathematical English.  Proofs
are omitted --- they are machine-checked by Lean.  The purpose is to let a
human verifier confirm that the formalisation captures the intended mathematics.
\end{abstract}

\tableofcontents
\bigskip

%===================================================================
\section{Overview}
%===================================================================

This file builds on \texttt{Containers.Basic} to equip containers and their
morphisms with the structure of a \emph{category} $\mathbf{Cont}$.

\begin{itemize}
  \item \textbf{Objects} are containers $C = (S \triangleleft P)$.
  \item \textbf{Morphisms} $(S \triangleleft P) \to (S' \triangleleft P')$ are
        container morphisms: a forward shape map $f : S \to S'$ together with a
        \emph{contravariant} position map $f^\sharp_s : P'(f(s)) \to P(s)$ for
        each $s$.
\end{itemize}

The key insight is that the contravariance of the position map makes all three
category laws hold \emph{definitionally} (by \lean{rfl}) --- no transports are
needed.  This contrasts sharply with the directed-container file, where genuine
dependent transports appear.

The file also establishes binary coproducts and products in $\mathbf{Cont}$,
completing the algebra of container constructors.

%===================================================================
\section{Identity and composition}
%===================================================================

% --- 1. ContainerMorphism.id ----------------------------------------
\subsection{Definition: \lean{ContainerMorphism.id}}

\begin{defn}[Identity morphism]
The \textbf{identity} container morphism on $C = (S \triangleleft P)$ is the
pair $(\mathrm{id}_S,\, \mathrm{id}_{P(s)})$: identity on shapes and identity
on positions.
\end{defn}

\paragraph{Lean.}
\begin{lstlisting}[style=lean]
def ContainerMorphism.id (C : Container) : ContainerMorphism C C where
  onShapes := id
  onPos := fun _ => id
\end{lstlisting}

\paragraph{Haskell.}
\begin{lstlisting}[style=haskell]
idMorph :: Morph s p s p
idMorph = Morph id (\_ -> id)
\end{lstlisting}


% --- 2. ContainerMorphism.comp --------------------------------------
\subsection{Definition: \lean{ContainerMorphism.comp}}

\begin{defn}[Composition of container morphisms]
Given morphisms $\varphi = (u, g) : C \to D$ and $\psi = (v, h) : D \to E$,
their \textbf{composite} $\psi \circ \varphi : C \to E$ is
\[
  (\psi \circ \varphi)_{\text{shapes}} \;=\; v \circ u, \qquad
  (\psi \circ \varphi)^\sharp_s \;=\; g_s \circ h_{u(s)}.
\]
Shape maps compose forwards; position maps compose \emph{backwards}, with
$\psi$'s pullback applied first, then $\varphi$'s.
\end{defn}

\paragraph{Lean.}
\begin{lstlisting}[style=lean]
def ContainerMorphism.comp {C D E : Container}
    (phi : ContainerMorphism C D) (psi : ContainerMorphism D E) :
    ContainerMorphism C E where
  onShapes := fun s => psi.onShapes (phi.onShapes s)
  onPos := fun s => phi.onPos s o. psi.onPos (phi.onShapes s)
\end{lstlisting}

\paragraph{Haskell.}
\begin{lstlisting}[style=haskell]
compMorph :: Morph s p t q -> Morph t q u r -> Morph s p u r
compMorph (Morph u g) (Morph v h) =
  Morph (v . u) (\s -> g s . h (u s))
\end{lstlisting}


%===================================================================
\section{Extensionality}
%===================================================================

% --- 3. ContainerMorphism.ext ----------------------------------------
\subsection{Theorem: \lean{ContainerMorphism.ext}}

\begin{thm}[Extensionality (HEq form)]
Two container morphisms $\varphi, \psi : C \to D$ are equal if their shape maps
are equal and their position maps are heterogeneously equal ($\lean{HEq}$):
\[
  \varphi.\mathrm{onShapes} = \psi.\mathrm{onShapes}
  \;\wedge\;
  \varphi.\mathrm{onPos} \approx \psi.\mathrm{onPos}
  \;\implies\;
  \varphi = \psi.
\]
The heterogeneous equality is necessary because the type of
$\varphi.\mathrm{onPos}$ depends on $\varphi.\mathrm{onShapes}$; after
substituting the shape equality, it collapses to ordinary equality.
\end{thm}

\paragraph{Lean.}
\begin{lstlisting}[style=lean]
theorem ContainerMorphism.ext {C D : Container} :
    forall {phi psi : ContainerMorphism C D},
      phi.onShapes = psi.onShapes -> HEq phi.onPos psi.onPos -> phi = psi
\end{lstlisting}

\paragraph{Haskell.}
\begin{lstlisting}[style=haskell]
-- In Haskell there is no dependent type on onPos, so ordinary
-- structural equality suffices: equal onShapes and equal onPos
-- implies equal Morph values.
\end{lstlisting}

\begin{fidelitynote}
The use of \lean{HEq} (heterogeneous equality) is forced by the dependent typing
of \lean{onPos}: its type mentions \lean{onShapes}, so two position maps that
live in \emph{different} function types cannot be compared with \lean{=} until
the shape equality is substituted.  This subtlety disappears in Haskell where
position maps are not dependent.
\end{fidelitynote}


% --- 4. ContainerMorphism.ext' -----------------------------------------
\subsection{Theorem: \lean{ContainerMorphism.ext'}}

\begin{thm}[Extensionality (pointwise form)]
A more convenient extensionality principle: given that the shape maps agree
($\lean{hs} : \varphi.\mathrm{onShapes} = \psi.\mathrm{onShapes}$), two morphisms
are equal whenever their position maps agree after transporting positions along
\lean{hs}.  When the shape maps agree \emph{definitionally}, the transport is
the identity and the hypothesis holds by \lean{rfl}.
\end{thm}

\paragraph{Lean.}
\begin{lstlisting}[style=lean]
theorem ContainerMorphism.ext' {C D : Container} :
    forall {phi psi : ContainerMorphism C D} (hs : phi.onShapes = psi.onShapes),
      (forall s p, phi.onPos s p = psi.onPos s (congrFun hs s |-> p)) -> phi = psi
\end{lstlisting}

\paragraph{Haskell.}
\begin{lstlisting}[style=haskell]
-- Again, in Haskell ordinary pointwise equality of the fields suffices.
-- (equal onShapes and pointwise-equal onPos) => equal Morph values.
\end{lstlisting}

\begin{fidelitynote}
The expression \lean{congrFun hs s} $\triangleright$ \lean{p} transports the position $p$ along the
shape equality: it coerces $p : P'(\varphi(s))$ to $P'(\psi(s))$ so that both
sides of the position-map equation live in the same type.  This is purely a
Lean book-keeping device with no Haskell analogue.
\end{fidelitynote}


%===================================================================
\section{Category laws}
%===================================================================

% --- 5. ContainerMorphism.id_comp ------------------------------------
\subsection{Theorem: \lean{ContainerMorphism.id\_comp}}

\begin{thm}[Left unit law]
For every morphism $\varphi : C \to D$,
\[
  \mathrm{id}_C \circ \varphi \;=\; \varphi.
\]
This holds definitionally (by \lean{rfl}).
\end{thm}

\paragraph{Lean.}
\begin{lstlisting}[style=lean]
theorem ContainerMorphism.id_comp {C D : Container}
    (phi : ContainerMorphism C D) :
    (ContainerMorphism.id C).comp phi = phi
\end{lstlisting}

\paragraph{Haskell.}
\begin{lstlisting}[style=haskell]
-- Left unit law:  compMorph idMorph phi = phi
\end{lstlisting}


% --- 6. ContainerMorphism.comp_id ------------------------------------
\subsection{Theorem: \lean{ContainerMorphism.comp\_id}}

\begin{thm}[Right unit law]
For every morphism $\varphi : C \to D$,
\[
  \varphi \circ \mathrm{id}_D \;=\; \varphi.
\]
This holds definitionally (by \lean{rfl}).
\end{thm}

\paragraph{Lean.}
\begin{lstlisting}[style=lean]
theorem ContainerMorphism.comp_id {C D : Container}
    (phi : ContainerMorphism C D) :
    phi.comp (ContainerMorphism.id D) = phi
\end{lstlisting}

\paragraph{Haskell.}
\begin{lstlisting}[style=haskell]
-- Right unit law:  compMorph phi idMorph = phi
\end{lstlisting}


% --- 7. ContainerMorphism.comp_assoc ----------------------------------
\subsection{Theorem: \lean{ContainerMorphism.comp\_assoc}}

\begin{thm}[Associativity]
For every composable triple $\varphi : C \to D$, $\psi : D \to E$,
$\chi : E \to F$,
\[
  (\chi \circ \psi) \circ \varphi \;=\; \chi \circ (\psi \circ \varphi).
\]
This holds definitionally (by \lean{rfl}).
\end{thm}

\paragraph{Lean.}
\begin{lstlisting}[style=lean]
theorem ContainerMorphism.comp_assoc {C D E F : Container}
    (phi : ContainerMorphism C D) (psi : ContainerMorphism D E)
    (chi : ContainerMorphism E F) :
    (phi.comp psi).comp chi = phi.comp (psi.comp chi)
\end{lstlisting}

\paragraph{Haskell.}
\begin{lstlisting}[style=haskell]
-- Associativity:
-- compMorph (compMorph phi psi) chi = compMorph phi (compMorph psi chi)
\end{lstlisting}


%===================================================================
\section{The \lean{Category} typeclass and the \lean{Cont} instance}
%===================================================================

% --- 8. Category typeclass -------------------------------------------
\subsection{Structure: \lean{Category}}

\begin{defn}[Category typeclass]
A minimal, Mathlib-free category typeclass on a type $\mathrm{Obj} :
\mathbf{Type}_u$.  It bundles:
\begin{itemize}
  \item a $\mathrm{Hom}$-family $\mathrm{Hom} : \mathrm{Obj} \to \mathrm{Obj}
        \to \mathbf{Type}_v$,
  \item identity morphisms $\mathrm{id}_X : \mathrm{Hom}(X, X)$,
  \item composition $\mathrm{comp} : \mathrm{Hom}(X,Y) \to \mathrm{Hom}(Y,Z)
        \to \mathrm{Hom}(X,Z)$ (diagrammatic order),
  \item proofs of left unit, right unit, and associativity.
\end{itemize}
This mirrors \lean{CategoryTheory.Category} from Mathlib but is self-contained.
\end{defn}

\paragraph{Lean.}
\begin{lstlisting}[style=lean]
class Category (Obj : Type u) where
  Hom : Obj -> Obj -> Type v
  id  : (X : Obj) -> Hom X X
  comp : {X Y Z : Obj} -> Hom X Y -> Hom Y Z -> Hom X Z
  id_comp : forall {X Y : Obj} (f : Hom X Y), comp (id X) f = f
  comp_id : forall {X Y : Obj} (f : Hom X Y), comp f (id Y) = f
  assoc : forall {W X Y Z : Obj} (f : Hom W X) (g : Hom X Y) (h : Hom Y Z),
    comp (comp f g) h = comp f (comp g h)
\end{lstlisting}

\paragraph{Haskell.}
\begin{lstlisting}[style=haskell]
class Category hom where
  identity :: hom a a
  compose  :: hom a b -> hom b c -> hom a c
  -- law: compose identity f = f
  -- law: compose f identity = f
  -- law: compose (compose f g) h = compose f (compose g h)
\end{lstlisting}


% --- 9. Cont instance -----------------------------------------------
\subsection{Instance: \lean{Cont}}

\begin{defn}[The category of containers]
\textbf{Cont} is the category whose objects are containers and whose morphisms
are container morphisms, with the identity and composition defined above.  All
three category laws hold definitionally.
\end{defn}

\paragraph{Lean.}
\begin{lstlisting}[style=lean]
instance Cont : Category Container where
  Hom    := ContainerMorphism
  id     := ContainerMorphism.id
  comp   := ContainerMorphism.comp
  id_comp := ContainerMorphism.id_comp
  comp_id := ContainerMorphism.comp_id
  assoc  := ContainerMorphism.comp_assoc
\end{lstlisting}

\paragraph{Haskell.}
\begin{lstlisting}[style=haskell]
-- Container morphisms form a category.
-- Objects: containers (s, p).
-- Morphisms: Morph s p t q with idMorph and compMorph.
-- All laws hold by computation.
\end{lstlisting}


%===================================================================
\section{Binary coproducts in \lean{Cont}}
%===================================================================

The coproduct of $C = (S \triangleleft P)$ and $D = (T \triangleleft Q)$ has shape
type $S + T$ and position type given case-wise: an $\mathrm{inl}(s)$-shape has
positions $P(s)$; an $\mathrm{inr}(t)$-shape has positions $Q(t)$.

% --- 10. Container.coprod -------------------------------------------
\subsection{Definition: \lean{Container.coprod}}

\begin{defn}[Coproduct container]
The \textbf{coproduct} $C \oplus D$ of containers $C = (S \triangleleft P)$ and
$D = (T \triangleleft Q)$ is the container
\[
  (S + T) \;\triangleleft\; [P, Q],
\]
where $[P, Q](\mathrm{inl}(s)) = P(s)$ and $[P, Q](\mathrm{inr}(t)) = Q(t)$.
\end{defn}

\paragraph{Lean.}
\begin{lstlisting}[style=lean]
def Container.coprod (C D : Container) : Container where
  Shape := C.Shape .+. D.Shape     -- Sum type: C.Shape + D.Shape
  Pos   := Sum.elim C.Pos D.Pos
\end{lstlisting}

\paragraph{Haskell.}
\begin{lstlisting}[style=haskell]
data CoprodShape s t = InlShape s | InrShape t

coprodPos :: (s -> p) -> (t -> q) -> CoprodShape s t -> Either p q
coprodPos posC _    (InlShape s) = Left  (posC s)
coprodPos _    posD (InrShape t) = Right (posD t)
-- or simply: Either s t with positions (p `Either` q) case-wise
\end{lstlisting}


% --- 11. Container.inl ----------------------------------------------
\subsection{Definition: \lean{Container.inl}}

\begin{defn}[Left injection]
The \textbf{left injection} $\iota_C : C \to C \oplus D$ has shape map
$\mathrm{inl} : S \to S + T$ and position map the identity (the positions of an
$\mathrm{inl}(s)$-shape in $C \oplus D$ are exactly $P(s)$).
\end{defn}

\paragraph{Lean.}
\begin{lstlisting}[style=lean]
def Container.inl (C D : Container) : ContainerMorphism C (C.coprod D) where
  onShapes := Sum.inl
  onPos    := fun _ => id
\end{lstlisting}

\paragraph{Haskell.}
\begin{lstlisting}[style=haskell]
inlMorph :: Morph s p (CoprodShape s t) (Either p q)
inlMorph = Morph InlShape (\_ -> Left)
-- onPos sends Left p back to p (identity on Left branch)
\end{lstlisting}


% --- 12. Container.inr ----------------------------------------------
\subsection{Definition: \lean{Container.inr}}

\begin{defn}[Right injection]
The \textbf{right injection} $\iota_D : D \to C \oplus D$ has shape map
$\mathrm{inr} : T \to S + T$ and position map the identity.
\end{defn}

\paragraph{Lean.}
\begin{lstlisting}[style=lean]
def Container.inr (C D : Container) : ContainerMorphism D (C.coprod D) where
  onShapes := Sum.inr
  onPos    := fun _ => id
\end{lstlisting}

\paragraph{Haskell.}
\begin{lstlisting}[style=haskell]
inrMorph :: Morph t q (CoprodShape s t) (Either p q)
inrMorph = Morph InrShape (\_ -> Right)
\end{lstlisting}


% --- 13. ContainerMorphism.desc ------------------------------------
\subsection{Definition: \lean{ContainerMorphism.desc}}

\begin{defn}[Copairing]
Given $\varphi : C \to E$ and $\psi : D \to E$, their \textbf{copairing}
$[\varphi, \psi] : C \oplus D \to E$ has shape map $[u_\varphi, u_\psi]$
(dispatch on $\mathrm{inl}$/$\mathrm{inr}$) and position map:
on an $\mathrm{inl}(s)$-shape use $\varphi^\sharp_s$; on an
$\mathrm{inr}(t)$-shape use $\psi^\sharp_t$.
\end{defn}

\paragraph{Lean.}
\begin{lstlisting}[style=lean]
def ContainerMorphism.desc {C D E : Container}
    (phi : ContainerMorphism C E) (psi : ContainerMorphism D E) :
    ContainerMorphism (C.coprod D) E where
  onShapes := Sum.elim phi.onShapes psi.onShapes
  onPos := fun s => match s with
    | Sum.inl s => phi.onPos s
    | Sum.inr s => psi.onPos s
\end{lstlisting}

\paragraph{Haskell.}
\begin{lstlisting}[style=haskell]
desc :: Morph s p u r -> Morph t q u r
     -> Morph (CoprodShape s t) (Either p q) u r
desc (Morph u g) (Morph v h) = Morph
  (\sh -> case sh of { InlShape s -> u s; InrShape t -> v t })
  (\sh -> case sh of { InlShape s -> g s; InrShape t -> h t })
\end{lstlisting}


% --- 14. ContainerMorphism.inl_desc ---------------------------------
\subsection{Theorem: \lean{ContainerMorphism.inl\_desc}}

\begin{thm}[Left computation rule for copairing]
\[
  \iota_C \circ [\varphi, \psi] \;=\; \varphi.
\]
Holds definitionally (\lean{rfl}).
\end{thm}

\paragraph{Lean.}
\begin{lstlisting}[style=lean]
theorem ContainerMorphism.inl_desc {C D E : Container}
    (phi : ContainerMorphism C E) (psi : ContainerMorphism D E) :
    (Container.inl C D).comp (phi.desc psi) = phi
\end{lstlisting}

\paragraph{Haskell.}
\begin{lstlisting}[style=haskell]
-- compMorph inlMorph (desc phi psi) = phi
\end{lstlisting}


% --- 15. ContainerMorphism.inr_desc ---------------------------------
\subsection{Theorem: \lean{ContainerMorphism.inr\_desc}}

\begin{thm}[Right computation rule for copairing]
\[
  \iota_D \circ [\varphi, \psi] \;=\; \psi.
\]
Holds definitionally (\lean{rfl}).
\end{thm}

\paragraph{Lean.}
\begin{lstlisting}[style=lean]
theorem ContainerMorphism.inr_desc {C D E : Container}
    (phi : ContainerMorphism C E) (psi : ContainerMorphism D E) :
    (Container.inr C D).comp (phi.desc psi) = psi
\end{lstlisting}

\paragraph{Haskell.}
\begin{lstlisting}[style=haskell]
-- compMorph inrMorph (desc phi psi) = psi
\end{lstlisting}


% --- 16. ContainerMorphism.desc_eta ---------------------------------
\subsection{Theorem: \lean{ContainerMorphism.desc\_eta}}

\begin{thm}[Universal property of coproduct ($\eta$-rule)]
Every morphism $\chi : C \oplus D \to E$ is the copairing of its composites
with the injections:
\[
  [\chi \circ \iota_C,\; \chi \circ \iota_D] \;=\; \chi.
\]
Together with the two computation rules, this is the full universal property:
$\mathrm{Hom}(C \oplus D, E) \cong \mathrm{Hom}(C, E) \times \mathrm{Hom}(D, E)$
naturally in $E$.
\end{thm}

\paragraph{Lean.}
\begin{lstlisting}[style=lean]
theorem ContainerMorphism.desc_eta {C D E : Container}
    (chi : ContainerMorphism (C.coprod D) E) :
    ((Container.inl C D).comp chi).desc
      ((Container.inr C D).comp chi) = chi
\end{lstlisting}

\paragraph{Haskell.}
\begin{lstlisting}[style=haskell]
-- desc (compMorph inlMorph chi) (compMorph inrMorph chi) = chi
\end{lstlisting}


% --- 17. ContainerMorphism.desc_unique ------------------------------
\subsection{Theorem: \lean{ContainerMorphism.desc\_unique}}

\begin{thm}[Uniqueness of the copairing mediating morphism]
For every cocone $(\varphi : C \to E,\; \psi : D \to E)$, there is a
\emph{unique} morphism $\chi : C \oplus D \to E$ with
$\chi \circ \iota_C = \varphi$ and $\chi \circ \iota_D = \psi$, namely
$\chi = [\varphi, \psi]$.
\end{thm}

\paragraph{Lean.}
\begin{lstlisting}[style=lean]
theorem ContainerMorphism.desc_unique {C D E : Container}
    (phi : ContainerMorphism C E) (psi : ContainerMorphism D E)
    (chi : ContainerMorphism (C.coprod D) E)
    (h1 : (Container.inl C D).comp chi = phi)
    (h2 : (Container.inr C D).comp chi = psi) :
    chi = phi.desc psi
\end{lstlisting}

\paragraph{Haskell.}
\begin{lstlisting}[style=haskell]
-- If compMorph inlMorph chi = phi and compMorph inrMorph chi = psi
-- then chi = desc phi psi
\end{lstlisting}


%===================================================================
\section{Binary products in \lean{Cont}}
%===================================================================

The product of $C = (S \triangleleft P)$ and $D = (T \triangleleft Q)$ has shape
type $S \times T$ and position type $P(s) + Q(t)$: a \emph{coproduct} of
position sets (not a product).  This is the categorical product of the
underlying polynomial functors.

% --- 18. Container.prod ---------------------------------------------
\subsection{Definition: \lean{Container.prod}}

\begin{defn}[Product container]
The \textbf{product} $C \times D$ of containers $C = (S \triangleleft P)$ and
$D = (T \triangleleft Q)$ is the container
\[
  (S \times T) \;\triangleleft\; \bigl((s, t) \mapsto P(s) + Q(t)\bigr).
\]
Note: the positions are the \emph{sum} (coproduct) $P(s) + Q(t)$, not the
product $P(s) \times Q(t)$.  This is the standard polynomial product.
\end{defn}

\paragraph{Lean.}
\begin{lstlisting}[style=lean]
def Container.prod (C D : Container) : Container where
  Shape := C.Shape × D.Shape
  Pos   := fun st => C.Pos st.1 .+. D.Pos st.2  -- coproduct of position types
\end{lstlisting}

\paragraph{Haskell.}
\begin{lstlisting}[style=haskell]
-- Product container: shapes are pairs, positions are sums.
prodPos :: (s -> p) -> (t -> q) -> (s, t) -> Either p q
prodPos posC posD (s, t) = Left (posC s)  -- or Right (posD t)
-- i.e. Pos (s, t) = Either (Pos_C s) (Pos_D t)
\end{lstlisting}

\begin{fidelitynote}
The product of polynomial functors (containers) uses a \emph{coproduct} of
position types because the polynomial functor associated to $C$ at a set $X$ is
$\sum_s X^{P(s)}$, and the product of two such functors is
$\sum_{(s,t)} X^{P(s) + Q(t)}$ (by the exponential law $X^{A+B} \cong X^A
\times X^B$).  This is distinct from the ``set-level'' product which would use
$P(s) \times Q(t)$ for positions.
\end{fidelitynote}


% --- 19. Container.fst -----------------------------------------------
\subsection{Definition: \lean{Container.fst}}

\begin{defn}[First projection]
The \textbf{first projection} $\pi_1 : C \times D \to C$ has shape map
$\mathrm{fst} : S \times T \to S$ and position map $\mathrm{inl} : P(s) \to
P(s) + Q(t)$ (the coproduct injection, going \emph{backwards} in position
variance).
\end{defn}

\paragraph{Lean.}
\begin{lstlisting}[style=lean]
def Container.fst (C D : Container) : ContainerMorphism (C.prod D) C where
  onShapes := Prod.fst
  onPos    := fun _ => Sum.inl
\end{lstlisting}

\paragraph{Haskell.}
\begin{lstlisting}[style=haskell]
fstMorph :: Morph (s, t) (Either p q) s p
fstMorph = Morph fst (\_ -> Left)
-- onPos sends (Left p) back to p; the backward map is Sum.inl
\end{lstlisting}


% --- 20. Container.snd -----------------------------------------------
\subsection{Definition: \lean{Container.snd}}

\begin{defn}[Second projection]
The \textbf{second projection} $\pi_2 : C \times D \to D$ has shape map
$\mathrm{snd} : S \times T \to T$ and position map $\mathrm{inr} : Q(t) \to
P(s) + Q(t)$.
\end{defn}

\paragraph{Lean.}
\begin{lstlisting}[style=lean]
def Container.snd (C D : Container) : ContainerMorphism (C.prod D) D where
  onShapes := Prod.snd
  onPos    := fun _ => Sum.inr
\end{lstlisting}

\paragraph{Haskell.}
\begin{lstlisting}[style=haskell]
sndMorph :: Morph (s, t) (Either p q) t q
sndMorph = Morph snd (\_ -> Right)
\end{lstlisting}


% --- 21. ContainerMorphism.pair ---------------------------------------
\subsection{Definition: \lean{ContainerMorphism.pair}}

\begin{defn}[Pairing]
Given $f : R \to C$ and $g : R \to D$, their \textbf{pairing}
$\langle f, g \rangle : R \to C \times D$ has
\begin{align*}
  \langle f, g \rangle_{\text{shapes}}(r)
    &\;=\; (f_{\text{shapes}}(r),\; g_{\text{shapes}}(r)), \\
  \langle f, g \rangle^\sharp_r
    &\;=\; [f^\sharp_r,\; g^\sharp_r]
      : P(f(r)) + Q(g(r)) \to R.\mathrm{Pos}(r).
\end{align*}
The position map case-splits via $[-, -]$ (eliminator for the sum).
\end{defn}

\paragraph{Lean.}
\begin{lstlisting}[style=lean]
def ContainerMorphism.pair {R C D : Container}
    (f : ContainerMorphism R C) (g : ContainerMorphism R D) :
    ContainerMorphism R (C.prod D) where
  onShapes := fun r => (f.onShapes r, g.onShapes r)
  onPos    := fun r => Sum.elim (f.onPos r) (g.onPos r)
\end{lstlisting}

\paragraph{Haskell.}
\begin{lstlisting}[style=haskell]
pairMorph :: Morph r rp s p -> Morph r rp t q
          -> Morph r rp (s, t) (Either p q)
pairMorph (Morph u g) (Morph v h) = Morph
  (\r -> (u r, v r))
  (\r -> either (g r) (h r))
\end{lstlisting}


% --- 22. ContainerMorphism.pair_fst -----------------------------------
\subsection{Theorem: \lean{ContainerMorphism.pair\_fst}}

\begin{thm}[Left computation rule for pairing]
\[
  \langle f, g \rangle \circ \pi_1 \;=\; f.
\]
Holds definitionally (\lean{rfl}).
\end{thm}

\paragraph{Lean.}
\begin{lstlisting}[style=lean]
theorem ContainerMorphism.pair_fst {R C D : Container}
    (f : ContainerMorphism R C) (g : ContainerMorphism R D) :
    (f.pair g).comp (Container.fst C D) = f
\end{lstlisting}

\paragraph{Haskell.}
\begin{lstlisting}[style=haskell]
-- compMorph (pairMorph f g) fstMorph = f
\end{lstlisting}


% --- 23. ContainerMorphism.pair_snd -----------------------------------
\subsection{Theorem: \lean{ContainerMorphism.pair\_snd}}

\begin{thm}[Right computation rule for pairing]
\[
  \langle f, g \rangle \circ \pi_2 \;=\; g.
\]
Holds definitionally (\lean{rfl}).
\end{thm}

\paragraph{Lean.}
\begin{lstlisting}[style=lean]
theorem ContainerMorphism.pair_snd {R C D : Container}
    (f : ContainerMorphism R C) (g : ContainerMorphism R D) :
    (f.pair g).comp (Container.snd C D) = g
\end{lstlisting}

\paragraph{Haskell.}
\begin{lstlisting}[style=haskell]
-- compMorph (pairMorph f g) sndMorph = g
\end{lstlisting}


% --- 24. ContainerMorphism.pair_eta -----------------------------------
\subsection{Theorem: \lean{ContainerMorphism.pair\_eta}}

\begin{thm}[Universal property of product ($\eta$-rule)]
Every morphism $\chi : R \to C \times D$ is the pairing of its composites
with the projections:
\[
  \langle \chi \circ \pi_1,\; \chi \circ \pi_2 \rangle \;=\; \chi.
\]
Together with the two computation rules, this is the full universal property:
$\mathrm{Hom}(R, C \times D) \cong \mathrm{Hom}(R, C) \times \mathrm{Hom}(R, D)$
naturally in $R$.
\end{thm}

\paragraph{Lean.}
\begin{lstlisting}[style=lean]
theorem ContainerMorphism.pair_eta {R C D : Container}
    (chi : ContainerMorphism R (C.prod D)) :
    (chi.comp (Container.fst C D)).pair
      (chi.comp (Container.snd C D)) = chi
\end{lstlisting}

\paragraph{Haskell.}
\begin{lstlisting}[style=haskell]
-- pairMorph (compMorph chi fstMorph) (compMorph chi sndMorph) = chi
\end{lstlisting}


% --- 25. ContainerMorphism.pair_unique --------------------------------
\subsection{Theorem: \lean{ContainerMorphism.pair\_unique}}

\begin{thm}[Uniqueness of the pairing mediating morphism]
For every cone $(f : R \to C,\; g : R \to D)$, there is a \emph{unique}
morphism $\chi : R \to C \times D$ with $\chi \circ \pi_1 = f$ and
$\chi \circ \pi_2 = g$, namely $\chi = \langle f, g \rangle$.
\end{thm}

\paragraph{Lean.}
\begin{lstlisting}[style=lean]
theorem ContainerMorphism.pair_unique {R C D : Container}
    (f : ContainerMorphism R C) (g : ContainerMorphism R D)
    (chi : ContainerMorphism R (C.prod D))
    (h1 : chi.comp (Container.fst C D) = f)
    (h2 : chi.comp (Container.snd C D) = g) :
    chi = f.pair g
\end{lstlisting}

\paragraph{Haskell.}
\begin{lstlisting}[style=haskell]
-- If compMorph chi fstMorph = f and compMorph chi sndMorph = g
-- then chi = pairMorph f g
\end{lstlisting}


%===================================================================
\section{Translation conventions}
%===================================================================

The following conventions supplement those in \texttt{Basic.tex}.

\begin{enumerate}
\item \textbf{Definitional equality.}  Many laws in this file hold
  \emph{definitionally} (the proof is \lean{rfl}).  This relies on Lean's
  $\eta$-rules for functions and structures.  In Haskell, equational laws are
  asserted as comments; they would be QuickCheck properties in practice.

\item \textbf{HEq and transport.}  The extensionality theorem
  \lean{ContainerMorphism.ext} uses heterogeneous equality (\lean{HEq}) because
  the position-map type depends on the shape map.  The pointwise variant
  \lean{ext'} uses an explicit transport ($\triangleright$).  Neither has a direct
  Haskell analogue.

\item \textbf{Product positions are sums.}  The product container has positions
  $P(s) + Q(t)$.  This is structurally dual to the coproduct (where shapes are
  sums and positions are case-wise), and reflects the polynomial-functor
  formula $X^{A+B} \cong X^A \times X^B$.

\item \textbf{Variance teaching.}  The injections $\iota_C, \iota_D$ are
  identity-on-positions; the projections $\pi_1, \pi_2$ carry coproduct
  injections backward on positions.  This duality captures the whole variance
  of containers at the level of (co)limits.
\end{enumerate}

\end{document}
